\section{怎样写日记}

我们注意到对于记叙、议论、说明诸基本写作能力的培养，加以必要的反复的练习，就可以逐步写好记叙文、议论文或说明文。但是这还不够，因为我们在学习、生活或工作中，大量需要的是由于写作目的不同，所写文章的性质形态也各不相同的应用文体。例如：日记、书信、读书笔记、观察笔记、计划、总结以及演讲稿、会议纪要、调查报告、电报稿、通知、请示、报告、科学小论文等等。对于这内容各别，形式多样的应用文体，我们虽然没有必要一一掌握它们的写法，但对日记、书信、计划、总结等最常用的应用文体则应非常熟悉，并且注意练习。只有这样，才能全面地提高我们的写作能力，达到学以致用的目的。

我们就先来谈谈日记和书信。日记是个人每天生活的真实记录，它记录着自己每天的经历和感受；它既是我们人生旅途中的脚印，也是鞭策自己前进的武器。书信则是个人与对方（亲友、同学、同志）互相联系、交流思想的一种工具；它既是自己的学习、生活、工作情况的介绍和汇报，也是自己的思想、见解、感情的表达和倾吐，希望得到对方的回答和反响。日记的读者是自己，书信的读者是对方。这就决定了日记或书信具有独特写法的特点。

我们怎样来写日记呢？

\textbf{1. 留心生活，注意观察，从我出发，广搜材料。}
一般来说，初写日记时遇到的困难往往是不知道该写什么。在一些同学看来，学生的生活无非是上学读书，回家吃饭，自习睡觉；在一些社会青年看来，生活也无非是上班下班，吃饭睡觉罢了。他们总觉得生活没有多少变化，难有新鲜内容，常写就乏味了；随之而来的，就是写的内容千篇一律，写来写去老一套，公式化。

其实，生活本身总是丰富多彩的，每天都会有新的经历，新的感受，关键在于你得做个“有心人”，留心生活，注意观察，认真思考，广加搜集，那么，你一定能找到可记、值得记的东西。被列宁称赞为“辩证法的奠基人之一”的古代希腊哲学家赫拉克利特说过：你不能两次走进同一条河。因为当你第二次走进这条河流时，原来的那条河流早就变化了。不容易看出变化的河流尚且如此，中学生在校学习各种课程，在校内外参加各项活动，接触很多人和事，这些人和事天天都处在变化发展之中，值得一写的内容实在太多。我们试以一天从早到晚的活动举例：清早起来练了长跑，可以写长跑的情况和感受，从中领悟到长跑的哪些要领和由长跑引起的对人生历程的联想；早饭时，广播里传来祖国各条战线的喜讯，播唱了发人深思的歌曲，报道了令人意外的消息等等，自己听了很受鼓舞、感动或震惊，可以写一写自己的心得体会或心理活动过程；上午听数学老师讲解题方法，联系自己的学习实际，可以写点收获，也可以描画一下老师讲课的神态，表达一下自己对老师敬爱的心情；听历史老师讲中国三大发明，可以写一写自己的民族自豪感，或者就是自己是黄帝的子孙该怎样做发一议议论；下午参加科技小组活动做了一个生物标本，受到老师的夸奖，非常高兴，可以把制作过程和体会写下来；放学回家路上，看到同学帮拉车工人推车上坡，自己没有上前，想起来感到羞愧。可以记下自己的心理活动过程，找出差距，晚上在灯下看科学幻想小说，联想到前几天看到的科技成就展览，也能写一点自己受到的触动或者自己有过的科学幻想……这些事情要说平常，都很平常，但只要选择一两个，就能写出内容丰富的日记。何况生活正像河流一样，每天每时都在发生变化，自己接触的人和事也在不断变化，即使是听同一位老师上语文课吧，今天同昨天所讲的内容不同，感受也会有所不同，怎能说一模一样，没什么好记的呢！

看来，说自己的生活平常，老一套，没什么可写的，并不符合事实。只要你睁大自己的双眼，开动自己的脑筋，善于把看似平常的小事、生活的细节中的深刻意义发掘出来，就能感到世上的万事万物都在可写之列，就能捕捉到平凡生活中的闪光点。

\textbf{2. 选好题材，抓住重点，中心突出地写一件事。}

既然只要做个有心人，就会发现有写不完、写不尽的素材，而一篇日记，篇幅短小，总不能包罗万象吧，那该怎么办呢？这就要学会对生活素材进行筛选，选择有典型意义的凡人小事，突破一点，进行记叙，力求达到“以小见大”、“小中显大”的境地。

不少同志初写日记，往往是事无巨细，全部记下，结果往往使日记成了一天生活的流水账。俗话说：伤其十指，不如断其一指。与其把一天经历的事情都来个蜻蜓点水，还不如选取一点写深写透。例如，有的同志经过记录一段“流水账”之后，慢慢学会了抓住重点题材，抛弃一般素材，像突出电影的“特写镜头”一样，只写使自己最感兴趣或最有感受的一件事，甚至还给每篇日记拟上一个标题，如《种树》、《我会修电灯了》、《看书法展览后》、《百米赛跑，起跑要快》、《读贺敬之的诗》、《少说空话，多做实事》、《替弟弟补功课所想到的》、《爸爸笑了》、《海边看日出》、《王老师一段语重心长的话》、《在图书馆的书架前》、《我同张庆的辩论》、《捉老鼠》、《考试碰上了难题》、《今天我虚度了吗？》，等等。写这样的日记，不仅对于日后查考回忆和提高思想认识有所帮助，而且直接地训练了自己观察生活、表达感受的写作能力。因为在考虑写什么的时候，就锻炼了立意、选材的能力。在考虑表达的时候，就锻炼了剪裁、布局、遣词造句的能力；又因为它的表达方式不拘，可以随兴之所至，经常变换记叙、描写、议论、抒情、说明的样式，这就锻炼了各种文体的写作能力；同时因为日记是每天都记的，也就是每天都在动脑、练笔，所以还锻炼了观察事物、分析问题以及经常整理思想的能力。

\textbf{3. 形式要活泼多样，不拘一格。}

写日记，除了标明每天的日期、天气以备查检回忆之外，一不要受表达方式的限制，二不必受篇幅的限制。哪一天见闻丰富，感受深刻，愿意尽情写，可以写成洋洋洒洒的千字长文；哪一天只有一点印象，一点感受，愿意简要写，可以写成精炼隽永的三言两语。想用记叙方式写一件小事，想用抒情方式抒发一下情怀，想用描写方式写一下所见的景色，想用说明方式介绍一下观察到的某种事物，想用议论方式对生活现象发一通感慨或表达一下自己的看法，可以随意选择，灵活运用。每篇日记的体裁也不必限得太死，可以是散文、随笔，可以是摘录、提要，可以是诗、警句，可以是杂感、小品。总之，不仅要使日记的内容丰富多彩，而且要使日记的形式灵活多样，使自己越写越想写，越看越爱看。

应该特别注意：

（1）一般说来，日记是写给自己看的，当然应该写自己的真情实感，而不可虚饰浮夸，无病呻吟。而且，这种真情实感必须是积极的，决不可格调低沉，自我欣赏，表现了不健康的情调，甚至用错误的立场、观点来看待社会，看待人生。我们既不要把日记写成检讨书、决心书之类，使日记成为和行动脱离的空话，更不可把日记写成发泄私愤甚至攻击现实的工具，而应利用写日记来督促、激励自己积极向上，奋勇前进。

（2）一定要有恒心和毅力，坚持不懈地写下去，而不要三天打鱼，两天晒网，甚至因为一忙，就把日记记的时间挤掉了。著名气象学家竺可桢数十年如一日，写下了四十几本日记，鲁迅先生从一九一二年五月五日起，到逝世的前两天的一九三六年十月十七日为止，写了近二十五年的日记。我们应把他们作为自己锻炼意志的榜样，既然开始记了，就要坚持不懈地记下去。毫无疑问，这对培养自己的高尚情操和提高自己的写作水平将有莫大的帮助。
\newpage